\section{Find the Minimum and Maximum in a BST}
\label{sec:bst-min-max}

\problemstatement{Given the root of a BST, find its minimum value and its
maximum value.}

\begin{patternbox}
\emph{``Minimum/maximum of a BST''} requires no comparison logic at all,
unlike Section~\ref{sec:bst-search} -- the BST invariant alone pins the
minimum and maximum to two fixed, unambiguous locations: the
\textbf{leftmost} node and the \textbf{rightmost} node.
\end{patternbox}

\subsection*{Understanding the problem carefully}

By the BST invariant, every left child is smaller than its parent, and
every right child is larger. So starting from the root and always moving
left strictly decreases the value at every step, until no left child
remains -- at which point we are standing on a value smaller than every
other value in the tree (every node we passed is larger, and every node in
their right subtrees, which we skipped, is even larger still). The
symmetric argument, always moving right, finds the maximum.

\begin{ideabox}[Walk to the Edge, No Comparisons Needed]
For the minimum: starting at the root, repeatedly move to the left child
until there is none; return the value at that point. For the maximum:
symmetrically, repeatedly move right until there is none.
\end{ideabox}

\subsection*{Why no value comparisons are needed during the walk}

This is a direct structural consequence of the BST invariant, not
something that needs separate verification at each step: the invariant
\emph{guarantees}, for every node visited, that its left child (if any) is
strictly smaller, so simply choosing ``left, if it exists'' at every step
is already provably optimal -- there is no need to compare the current
node's value to anything, since the tree's shape alone already encodes
the answer's location.

\subsection*{Algorithm}

\begin{enumerate}
  \item \textsc{FindMin}: starting at \textit{root}, while the current
        node has a left child, move to it. Return the current node's
        value.
  \item \textsc{FindMax}: symmetrically, while the current node has a
        right child, move to it. Return the current node's value.
\end{enumerate}

\subsection*{Dry run}

\dryrun{Take the BST with root $5$, left child $3$ (with left child $1$
and right child $4$), right child $8$ (with right child $9$).}

\textbf{Min}: $5 \to 3 \to 1$ (no left child). Minimum is $1$.

\textbf{Max}: $5 \to 8 \to 9$ (no right child). Maximum is $9$.

\subsection*{C++ implementation}

\begin{lstlisting}
#include <bits/stdc++.h>
using namespace std;

struct TreeNode {
    int val;
    TreeNode* left;
    TreeNode* right;
    TreeNode(int v) : val(v), left(nullptr), right(nullptr) {}
};

int findMin(TreeNode* root) {
    while (root->left) root = root->left;
    return root->val;
}

int findMax(TreeNode* root) {
    while (root->right) root = root->right;
    return root->val;
}

int main() {
    TreeNode* root = new TreeNode(5);
    root->left = new TreeNode(3);
    root->right = new TreeNode(8);
    root->left->left = new TreeNode(1);
    root->left->right = new TreeNode(4);
    root->right->right = new TreeNode(9);

    cout << "Min: " << findMin(root) << ", Max: " << findMax(root) << endl;
    return 0;
}
\end{lstlisting}

\begin{complexitybox}
\textbf{Time Complexity.} $O(h)$, where $h$ is the tree's height.
\textbf{Space Complexity.} $O(1)$ -- the iterative walk needs no
recursion stack at all, unlike most other BST operations in this chapter.
\end{complexitybox}

\begin{takeawaybox}
Finding the minimum and maximum of a BST is the simplest possible
operation in this chapter precisely because it needs zero comparisons --
the structure alone answers the question. This ``walk to the leftmost /
rightmost node'' primitive is reused directly inside Delete
(Section~\ref{sec:bst-delete}) to find a replacement node, and inside
Predecessor/Successor (Section~\ref{sec:bst-pred-succ}).
\end{takeawaybox}
